\subsection{Real-time Video Generation}
\label{sec:real_time_gen}

Current video generative methods often require significant computational resources and considerable time to produce even short video clips. Systems relying on high-end hardware remain slow, often taking several minutes to generate just a few seconds of video. While the visual quality of these generated videos has seen substantial improvements in recent years, the extended processing times present a major bottleneck in practical applications and iterative design workflows. In fields where rapid prototyping and real-time adjustments are crucial—such as interactive entertainment, virtual reality, and video production—these long generation times limit creative flexibility and slow down the process of refining visual content.

Moreover, the slow speed of current methods limits their use in dynamic environments where real-time feedback is essential. In contexts like live-streaming or gaming, there’s an increasing demand for systems that can instantly render synthetic scenes and characters in response to user actions or changing inputs. Any delay in generating high-quality video clips can result in a suboptimal user experience, reducing immersion and interactivity. By focusing on real-time generation, \method seeks to enhance both the performance and applicability of video generation techniques, broadening their potential in fast-paced, interactive environments. Fig.~\ref{fig:streamer1} and \ref{fig:streamer2} provide such cases of real-time long video stream generation. We then introduce how to achieve this using our \method models.

\begin{figure*}[ht]
    \centering
    \includegraphics[width=1\linewidth]{figures/realtime/stream1.pdf}
    \caption{Generated 15 minutes video with 8 A100 GPUs in real-time 8 FPS.}
    \label{fig:streamer1}
\end{figure*}

\begin{figure*}[ht]
    \centering
    \includegraphics[width=1.\linewidth]{figures/realtime/stream2.pdf}
    \caption{Generated videos with a single RTX 4090 GPU, int8 and TensorRT quantization in real-time 20 FPS.}
    \label{fig:streamer2}
\end{figure*}

\subsubsection{Method}
To construct the real-time generation pipeline, we build upon a previously pre-trained version of the \method model. This design choice offers two primary benefits. First, starting from a well-trained model significantly accelerates convergence and improves training stability, as the model already possesses a strong initialization. Second, the pre-trained model has already captured valuable knowledge about motion patterns and continuous temporal dynamics, which can be directly inherited by the new pipeline to enhance the smoothness and coherence of generated content. Therefore, we adopt the pre-trained \method model as the foundation for our real-time pipeline~\footnote{Part of this section is introduced in a previously released tech report \emph{The Matrix}~\citep{feng2024matrixinfinitehorizonworldgeneration}.}.

The pre-trained \method model is designed to generate fixed-length video clips, typically 5 to 10 seconds long. To transform this into a real-time streaming model capable of generating continuous video streams with no predefined length (potentially infinite), two key modifications are required:
\begin{itemize}
    \item \textbf{Streaming pipeline adaptation.} We replace the static generation process, which produces a fixed-length sequence of video tokens in a single pass, with a streaming mechanism that incrementally generates video tokens. In this streaming setup, video tokens are processed through a denoising queue — each time the oldest token is generated and dequeued, a new token is added to the bottom of the queue, enabling continuous generation without length constraints.
    \item \textbf{Real-time acceleration.} In addition to enabling streaming, we optimize the generation speed to ensure the system meets real-time performance requirements, \ie the pipeline must generate new frames fast enough to keep up with real-time playback.
\end{itemize}
In the following sections, we will describe each of these modifications in detail, including their design rationale, technical implementation, and empirical performance.


\subsubsection{Streaming Video Generation}
Traditional Diffusion Transformer (DiT)~\citep{dit} models are constrained in their ability to generate videos longer than a few seconds, even when employing significant spatial and temporal compression through Variational Autoencoders (VAEs)~\citep{kingma2013auto}. This limitation stems primarily from the computational and memory-intensive nature of attention mechanisms over extended temporal sequences. To overcome this challenge, we introduce Streamer, a novel approach that leverages a sliding temporal window to manage temporal dependencies efficiently, enabling the generation of long or even infinite-length videos in a streaming fashion.


\paragraph{Shift window denoise process models.} The key innovation of Streamer lies in its assumption that temporal dependencies are confined within a limited time window. By focusing attention computations within this window, Streamer significantly reduces computational overhead while maintaining video continuity. Specifically, Streamer processes video tokens in a queue, denoising them over multiple noise levels within the window. After a fixed number of denoising steps, the leftmost token (with the lowest noise level) is dequeued and cached, while a new token with Gaussian noise is appended to the rightmost position. This sliding mechanism ensures that the window maintains a consistent length while facilitating the generation of continuous video frames.

\paragraph{Training and inference.}
Streamer is fine-tuned from a pre-trained DiT model. During training, a sequence of $2w$ video tokens is sampled, where $w$ is the size of the sliding window (typically set to $w=T$, the number of diffusion solver steps). The first $w$ tokens are used solely for `warming up' the model and do not contribute to the loss computation. The loss is calculated only on the last $w$ tokens, ensuring that the model learns to generate coherent video frames within the window. At inference time, the same warmup strategy is employed: the first $w$ tokens are discarded, and the generated video begins from the $(w+1)$-th token.

\paragraph{Ensuring continuity during training and inference.}
To maintain smooth transitions between consecutive windows, cached tokens are re-introduced into the token queue at a noise level of 0. This allows previously generated tokens to participate in the denoising process of subsequent windows, ensuring temporal consistency and continuity in the generated video.

Streamer represents a significant advancement in video generation by adapting DiT models for streaming applications. Its sliding temporal window approach addresses the limitations of traditional DiT models, enabling the efficient generation of infinite-length videos while maintaining high quality and temporal coherence. This method opens new possibilities for real-time video generation and long-form content creation. We summarize its advantages as follows:
\begin{itemize}
    \item Infinite-Length Video Generation: By leveraging a sliding window mechanism, Streamer can generate videos of arbitrary length without incurring prohibitive computational costs;
\item Efficient Attention Computation: Limiting attention computations to a confined temporal window reduces memory and processing demands;
\item Seamless Continuity: The caching and re-introduction of tokens ensure smooth transitions between windows, preserving video quality over extended durations.
\end{itemize}


\subsubsection{Consistency Model Distillation}
After successfully extending the DiT model to Streamer for infinite-length video generation, the next critical step is to address the challenge of achieving real-time rendering of the simulated world. While Streamer enables the generation of long or even infinite videos by efficiently managing temporal dependencies within a sliding window, the computational demands of the diffusion process still pose a bottleneck for real-time applications. To overcome this, we propose a novel integration of Streamer with Consistency Models, a cutting-edge approach for accelerating diffusion-based generative processes. Specifically, we leverage the Latent Consistency Model (LCM~\citep{luo2023latent}) and its video version (VideoLCM~\citep{wang2023videolcm}), which distill the original diffusion process and  class-free guidance into a highly efficient four-step consistency model. This distillation process significantly reduces the number of sampling steps required for high-quality generation while maintaining the temporal coherence and streaming capabilities inherent to Streamer.

The integration of LCM with Streamer is designed to preserve the denoising window mechanism, ensuring that the sliding temporal window continues to manage dependencies effectively. During training, we optimize this combined framework to balance efficiency and quality. The result is a 10 - 20 $\times$ acceleration in inference speed, enabling the model to achieve a rendering rate of 8 - 16 FPS. This dramatic improvement in performance makes the system viable for real-time applications, such as interactive simulations, live video synthesis, and dynamic content generation in virtual environments. By combining the strengths of Streamer—its ability to generate infinite-length videos with temporal coherence—and LCM—its efficiency in accelerating the diffusion process—we bridge the gap between high-quality video generation and real-time responsiveness. This advancement not only enhances the practicality of diffusion models for streaming applications but also opens up new possibilities for immersive, interactive, and real-time media creation.

\paragraph{Quantization for running in customer-level devices.} While the generating speed has been accelerated to real-time levels, deploying the model on consumer-level devices remains challenging due to the high computational and memory requirements, even for high-end GPUs like the NVIDIA 4090. To address this, we introduce quantization techniques to optimize the model for efficient deployment. Specifically, we employ two distinct quantization strategies: int8 quantization~\citep{torchao} for attention layers and linear heads, and TensorRT quantization~\citep{tensorrt} for the overall model. The int8 quantization approach significantly reduces memory consumption by converting weights and activations to 8-bit integers, while maintaining the overall generation quality. However, this method provides minimal acceleration in generating speed, as it primarily focuses on memory efficiency rather than computational optimization. On the other hand, TensorRT quantization offers a more comprehensive solution, enabling substantial acceleration in generation speed. This technique allows the model to achieve real-time performance of 8 FPS even on a single 4090 GPU, making it feasible for consumer-level devices. However, TensorRT quantization introduces a modest increase in network error, which can be observed even when its built-in error-checking mechanisms are enabled. This error manifests as slight deviations in output quality, such as minor artifacts or inconsistencies in the generated video. In some cases, it may also lead to measurable instability, such as flickering or temporal incoherence. Despite these trade-offs, the combination of int8 and TensorRT quantization provides a balanced approach to optimizing the model for real-time, consumer-level deployment, ensuring both efficiency and practicality while maintaining acceptable generation quality. By carefully tuning the quantization parameters and leveraging TensorRT’s error-checking features, we mitigate these issues to the greatest extent possible, enabling robust and efficient performance on consumer-grade hardware.
